\documentclass[11pt,a4paper]{article}
\usepackage[margin=1in]{geometry}
\usepackage{amsmath,amssymb,amsthm}
\usepackage{algorithm}
\usepackage{algpseudocode}
\usepackage{booktabs}
\usepackage{hyperref}

\theoremstyle{definition}
\newtheorem{definition}{\textbf{Definition}}[section]
\newtheorem{theorem}{\textbf{Theorem}}[section]
\newtheorem{lemma}[theorem]{\textbf{Lemma}}
\newtheorem{remark}{\textbf{Remark}}[section]

\title{\textbf{Monte Carlo Tree Search with Learned Networks (AlphaZero, MuZero)}}
\author{\textbf{Team Scaibu} \\ \small Email: \texttt{jaydeep.9664920749@gmail.com} \\ \small \textbf{Architecture and Core Engine Team}}
\date{\textbf{October 2026}}

\begin{document}
\maketitle

\section{Definitions and Cognitive Mental Models}

\subsection{Formal Mathematical Definition}
\textbf{Definition (Predictor-Guided Tree Search and Policy Iteration):}
Classical Monte Carlo Tree Search (MCTS) relies on random rollouts to estimate terminal state values, suffering from exponential sample complexity in deep game trees. \textbf{AlphaZero and MuZero} replace random rollouts with dual-headed deep neural networks that evaluate state value $v_\theta(s)$ and output prior action probabilities $P_\theta(a | s)$.

At each search step from state $s$, action selection follows the \textbf{Polynomial Upper Confidence Trees (PUCT)} algorithm:
{\boldmath
\begin{equation}
a^* = \arg\max_{a \in \mathcal{A}} \left[ Q(s, a) + U(s, a) \right]
\end{equation}
}
where the exploration term $U(s, a)$ is dynamically scaled by:
{\boldmath
\begin{equation}
U(s, a) = c(s) \cdot P_\theta(s, a) \cdot \frac{\sqrt{\sum_b N(s, b)}}{1 + N(s, a)}
\end{equation}
}
and the exploration rate grows logarithmically with total visits:
{\boldmath
\begin{equation}
c(s) = c_{\text{init}} + \ln\left( \frac{\sum_b N(s, b) + c_{\text{base}} + 1}{c_{\text{base}}} \right)
\end{equation}
}
After executing $M$ tree simulations, search visit counts $\{N(s, a)\}$ construct an improved policy distribution $\pi(a | s)$:
{\boldmath
\begin{equation}
\pi(a | s) = \frac{N(s, a)^{1/\tau}}{\sum_{b} N(s, b)^{1/\tau}}
\end{equation}
}
where $\tau > 0$ is the search temperature. The neural network is trained via self-play to minimize cross-entropy against $\pi$ and mean squared error against game outcome $z$:
{\boldmath
\begin{equation}
\mathcal{L}(\theta) = (z - v_\theta(s))^2 - \sum_{a} \pi(a | s) \ln P_\theta(a | s) + c_{\text{reg}}\|\theta\|^2
\end{equation}
}

\subsection{Expanded Non-Technical Definition (Intuition for Anyone)}
\textbf{Intuitive Conceptual Definition:}
\begin{enumerate}
    \item \textbf{Combining Gut Feeling with Deep Calculation}: A grandmaster in chess glances at the board and immediately ``feels'' 2 or 3 promising moves (neural network prior). Then, they calculate several moves ahead along those exact lines (tree search). AlphaZero combines these two modes of thinking perfectly.
    \item \textbf{The PUCT Balance}: At the start of thinking, the agent tries moves that look promising according to its gut intuition ($P(s, a)$). As it simulates more games, it trusts actual measured outcomes ($Q(s, a)$), guaranteeing that bad traps are discovered and discarded.
    \item \textbf{MuZero (Learning the Rules from Scratch)}: AlphaZero required knowing the official rules of chess or Go. MuZero goes further: it learns its own internal video game engine, allowing it to imagine future moves without ever being told how the game works!
    \item \textbf{Bootstrapping to Superhuman Mastery}: The tree search produces decisions that are strictly stronger than the raw neural network. By training the neural network to predict its own tree search results, the system pulls itself up by its own bootstraps into superhuman intelligence.
\end{enumerate}

\subsection{Child-Friendly Mental Model (Physical Metaphor)}
Imagine exploring a giant hedge maze to find a treasure chest:
\begin{itemize}
    \item \textbf{The Random Wanderer}: Runs down random hallways, gets lost, and takes a million years to find the chest!
    \item \textbf{The Magic Compass (Neural Prior)}: An enchanted compass points in the general direction of the treasure, suggesting which 2 hallways look best.
    \item \textbf{The Scout Team (MCTS)}: You send 50 fast scouts down those promising paths. Scouts who find dead ends report back ($Q(s, a)$ falls). Scouts who find golden corridors report back ($Q(s, a)$ rises).
    \item \textbf{Confident Stride}: You follow the path where the most scouts returned happily ($N(s, a)$ is highest)!
\end{itemize}

\section{Core Mathematical Equations and Definitions}

\subsection{PUCT Exploration Term}
{\boldmath
\begin{equation}
U(s, a) = \left( c_{\text{init}} + \ln\left( \frac{\sum_b N(s, b) + c_{\text{base}} + 1}{c_{\text{base}}} \right) \right) P(s, a) \frac{\sqrt{\sum_b N(s, b)}}{1 + N(s, a)}
\end{equation}
}
\begin{itemize}
    \item \textbf{Formal Definition}: The upper confidence bound exploration bonus balancing network prior against search visit density.
    \item \textbf{What It Means}: Guarantees unvisited promising actions are explored early, but asymptotic selection is dominated by true value $Q$.
    \item \textbf{Why It Is Important}: Prevents search from ignoring valid tactical moves that the initial neural prior undervalued.
\end{itemize}

\subsection{Temperature-Scaled Visit Policy}
{\boldmath
\begin{equation}
\pi(a | s) = \frac{N(s, a)^{1/\tau}}{\sum_b N(s, b)^{1/\tau}}
\end{equation}
}
\begin{itemize}
    \item \textbf{Formal Definition}: Normalized Boltzmann visit distribution parameterized by exploration temperature $\tau$.
    \item \textbf{What It Means}: When $\tau = 1$, action probability equals empirical visit proportion; as $\tau \to 0$, it converges to the greedy argmax action.
    \item \textbf{Why It Is Important}: Provides soft targets for self-play policy training while allowing deterministic play at tournament evaluation.
\end{itemize}

\subsection{Root Value Expectation}
{\boldmath
\begin{equation}
v_{\text{root}} = \sum_{a \in \mathcal{A}} \frac{N(s, a)}{\sum_b N(s, b)} Q(s, a)
\end{equation}
}
\begin{itemize}
    \item \textbf{Formal Definition}: Empirical visit-weighted mean of subtree action values at the root node.
    \item \textbf{What It Means}: Reflects the aggregate evaluation of the state incorporating all deep tactical simulations.
    \item \textbf{Why It Is Important}: Serves as the high-accuracy regression target for training the neural network's scalar value head.
\end{itemize}

\section{Rigorous Mathematical Analysis Checklist}

\begin{itemize}
    \item \textbf{Notation}: $N(s, a)$ (visit counts), $Q(s, a)$ (mean action value), $P(s, a)$ (network prior), $\tau$ (search temperature), $c_{\text{puct}}$ (exploration constant).
    \item \textbf{Epistemic Status}: Formulated by Silver et al. (DeepMind, 2017--2020 across AlphaGo, AlphaZero, MuZero), crowning triumph of deep search.
    \item \textbf{Dependency Graph}: State $s \to (\mathbf{P}, v) \to \text{PUCT Selection} \to \text{Expansion \& Backup} \to \{N(s, a)\} \to \boldsymbol{\pi} \to \text{Action}$.
    \item \textbf{Dimensions}: Discrete action space of size $|\mathcal{A}| = A$, simulation count $M$.
    \item \textbf{Regimes}: Combinatorial games (Go, Chess, Shogi), LLM multi-step reasoning, mathematical theorem proving, chip floorplanning.
    \item \textbf{Invariants}: If all visit counts $N(s, a)$ are identical and $\tau = 1$, target policy $\boldsymbol{\pi}$ is exactly uniform $1/A$.
    \item \textbf{Isomorphisms}: PUCT is an asymmetric Bayesian bandit problem with Dirichlet priors and logarithmic exploration bonuses.
\end{itemize}

\section{Theoretical Theorems and Formal Proofs}

\begin{theorem}[Asymptotic Convergence of PUCT to Optimal Minimax Action]
For any finite discrete Markov decision process with bounded rewards, as the number of MCTS simulations $M = \sum_a N(s, a) \to \infty$, the probability of selecting an action with suboptimal expected return converges to zero.
\end{theorem}

\begin{proof}
Let $a^* = \arg\max_{a} Q^*(s, a)$ be the unique optimal action, and let $a'$ be any suboptimal action such that $\Delta = Q^*(s, a^*) - Q^*(s, a') > 0$.
The PUCT selection condition requires choosing $a'$ over $a^*$ only if:
{\boldmath
\begin{equation}
Q(s, a') + U(s, a') \ge Q(s, a^*) + U(s, a^*)
\end{equation}
}
Notice that the exploration term $U(s, a)$ satisfies:
{\boldmath
\begin{equation}
U(s, a') = c(s) P(s, a') \frac{\sqrt{M}}{1 + N(s, a')}
\end{equation}
}
As $M \to \infty$, by the law of large numbers on tree rollouts, the sample mean estimates converge almost surely to their true values: $Q(s, a) \xrightarrow{\text{a.s.}} Q^*(s, a)$.
Suppose for contradiction that $N(s, a') \to \infty$ at rate comparable to $M$, such that $N(s, a') / M \ge \delta > 0$.
Then:
{\boldmath
\begin{equation}
U(s, a') = \mathcal{O}\left( \frac{\sqrt{M}}{\delta M} \right) = \mathcal{O}\left( \frac{1}{\sqrt{M}} \right) \to 0
\end{equation}
}
Consequently, for large $M$:
{\boldmath
\begin{equation}
Q(s, a') + U(s, a') \to Q^*(s, a') < Q^*(s, a^*) \le Q(s, a^*) + U(s, a^*)
\end{equation}
}
This violates the selection inequality. Therefore, $a'$ cannot be chosen infinitely often, and its visit count must satisfy $N(s, a') = \mathcal{O}(\sqrt{M})$.
Hence:
{\boldmath
\begin{equation}
\lim_{M \to \infty} \frac{N(s, a^*)}{M} = 1
\end{equation}
}
Under temperature $\tau \le 1$, $\pi(a^* | s) \to 1$ and $\pi(a' | s) \to 0$, proving asymptotic convergence to the optimal action.
\end{proof}

\begin{theorem}[Invariance of Uniform Visit Policy Under Temperature Scaling]
If the search visit counts are uniformly distributed $N(s, a) = K$ for all $a \in \mathcal{A}$, then the target policy $\pi(a | s) \equiv \frac{1}{|\mathcal{A}|}$ for all valid temperature parameters $\tau > 0$.
\end{theorem}

\begin{proof}
Let $|\mathcal{A}| = A$. Suppose $N(s, a) = K$ for all $a \in \{1, \dots, A\}$ with $K > 0$.
By definition of the visit policy:
{\boldmath
\begin{equation}
\pi(a | s) = \frac{N(s, a)^{1/\tau}}{\sum_{b=1}^A N(s, b)^{1/\tau}} = \frac{K^{1/\tau}}{\sum_{b=1}^A K^{1/\tau}}
\end{equation}
}
Factoring out the numerator:
{\boldmath
\begin{equation}
\pi(a | s) = \frac{K^{1/\tau}}{A \cdot K^{1/\tau}} = \frac{1}{A}
\end{equation}
}
Because $K^{1/\tau} > 0$ cancels identically, the resulting policy is strictly independent of temperature $\tau$.
\end{proof}

\section{Foundational Architectural Questions and Epistemic Meta-Questions}

\subsection{Architectural Question 1: Simulation Budget vs Inference Latency}
\textbf{Question:} How is the MCTS simulation budget determined in real-time agent deployments?\\
\textbf{Answer:} Tree search scales compute at test time: each simulation requires evaluating the neural network, adding latency linearly. In latency-sensitive settings (e.g. conversational agents, fast arcade games), budgets are set to $M \in [50, 200]$ simulations. In deep planning, code synthesis, or combinatorial logistics, budgets of $M \in [800, 5000]$ are standard to ensure exhaustive tactical coverage.

\textbf{Meta-Question:} Why is Dirichlet noise added at the root during self-play training?\\
\textbf{Answer:} Deep neural networks can assign zero probability prior to certain non-intuitive moves. Without noise, the PUCT search might never explore those actions, locking the self-play engine in an echo chamber of conventional opening moves. Adding Dirichlet noise $\operatorname{Dir}(\alpha)$ to the root prior forces exploration across all available branches.

\subsection{Architectural Question 2: MuZero Dynamics vs Environment Simulators}
\textbf{Question:} Why does MuZero outperform AlphaZero on arbitrary Atari environments?\\
\textbf{Answer:} AlphaZero requires an exact, perfect physical rules engine (the chess board logic) to advance the tree. In real-world environments and Atari games, perfect rules engines do not exist or are slow. MuZero trains an internal recurrent transition network that predicts next latent states and rewards, enabling search across arbitrary unmodeled games.

\textbf{Meta-Question:} How is MCTS applied to reasoning LLMs (e.g., OpenAI o1 style)?\\
\textbf{Answer:} Rather than branching over individual tokens, MCTS for reasoning LLMs treats intermediate reasoning paragraphs or calculation steps as actions. A Process Reward Model (PRM) or verifier acts as the state evaluation function, enabling tree search over multi-step mathematical proofs and logic chains.

\section{Algorithmic Implementation Specification}

\begin{algorithm}[H]
\caption{MCTS PUCT Action Selection and Target Extraction}
\begin{algorithmic}[1]
\Require Prior probabilities $\mathbf{P}$, visit counts $\mathbf{N}$, Q-values $\mathbf{Q}$, $c_{\text{init}}, c_{\text{base}}, \tau$
\Ensure Selected action $a^*$, PUCT scores $\mathbf{S}$, exploration terms $\mathbf{U}$, target policy $\boldsymbol{\pi}$, root value $v_{\text{root}}$
\State $total\_visits \gets \sum_{a=0}^{A-1} \mathbf{N}[a]$
\State $c_{\text{scale}} \gets c_{\text{init}} + \ln((total\_visits + c_{\text{base}} + 1.0) / c_{\text{base}})$
\State $best\_score \gets -\infty, \quad a^* \gets 0$
\State Allocate $\mathbf{S} \in \mathbb{R}^A, \quad \mathbf{U} \in \mathbb{R}^A, \quad \boldsymbol{\pi} \in \mathbb{R}^A$
\For{$a = 0$ \textbf{to} $A-1$}
    \State $u_a \gets c_{\text{scale}} \cdot \mathbf{P}[a] \cdot (\sqrt{total\_visits} / (1.0 + \mathbf{N}[a]))$
    \State $\mathbf{U}[a] \gets u_a$
    \State $\mathbf{S}[a] \gets \mathbf{Q}[a] + u_a$
    \If{$\mathbf{S}[a] > best\_score$}
        \State $best\_score \gets \mathbf{S}[a], \quad a^* \gets a$
    \EndIf
\EndFor
\If{$total\_visits == 0$}
    \State $\boldsymbol{\pi} \gets [1/A, \dots, 1/A], \quad v_{\text{root}} \gets \frac{1}{A}\sum_{a=0}^{A-1} \mathbf{Q}[a]$
\Else
    \State $sum\_pow \gets \sum_{a=0}^{A-1} (\mathbf{N}[a])^{1/\tau}$
    \For{$a = 0$ \textbf{to} $A-1$}
        \State $\boldsymbol{\pi}[a] \gets (\mathbf{N}[a])^{1/\tau} / sum\_pow$
    \EndFor
    \State $v_{\text{root}} \gets \sum_{a=0}^{A-1} (\mathbf{N}[a] / total\_visits) \cdot \mathbf{Q}[a]$
\EndIf
\State \Return $(a^*, \mathbf{S}, \mathbf{U}, \boldsymbol{\pi}, v_{\text{root}})$
\end{algorithmic}
\end{algorithm}

\section{Big-$\Theta$ Complexity Analysis}

\begin{itemize}
    \item \textbf{PUCT Selection Step}: $\Theta(A)$ operations over legal action space.
    \item \textbf{Policy Extraction Time}: $\Theta(A)$ power and normalization operations.
    \item \textbf{Root Value Computation}: $\Theta(A)$ multiplications and additions.
    \item \textbf{Space Complexity}: $\Theta(A)$ buffer for action vectors.
    \item \textbf{Tree Search Depth}: $\mathcal{D} = \Theta(H)$ where $H$ is the tree search depth limit.
\end{itemize}

\section{Single Canonical Repository Reference}

The complete deterministic scalar implementation, verification suite, and unit test benchmarks are published at:
\begin{center}
\url{https://github.com/Chief-Strategist-J/llm-obs-infra}
\end{center}

\end{document}
